

Low-latency applications, such as cloud gaming, cloud \ac{XR}, and the tactile internet, are poised to transform the future of communications and revolutionize industries including entertainment, healthcare, and education. These applications enable groundbreaking use cases such as real-time collaboration, lifelike simulations, and immersive virtual environments. However, their success is intrinsically tied to the performance of high-speed, low-latency networks. Variability in network conditions—such as fluctuating bandwidth, latency, jitter, or packet loss—can severely impact the \ac{QoE}, leading to performance failures in environments like cellular and Wi-Fi networks. For network operators, addressing these challenges is both a technical necessity and a business imperative, as it involves detecting and diagnosing the root causes of performance issues to maintain service reliability.

This thesis addresses these challenges by contributing insights and innovative solutions for detecting anomalies and diagnosing root causes in low-latency applications. The research begins with the foundational task of data collection, described in Chapter \ref{chap:measurements}. To create datasets representative of real-world LL application performance, we established multiple testbeds. First, we collected transmission opportunities (txops) from Orange's 4G network to simulate realistic network conditions. These txops were used to collect time series Key Performance Indicators (KPIs) from three commercial CG platforms. Additionally, a testbed was designed to collect Cloud VR KPIs under controlled Wi-Fi network conditions, using NVIDIA Cloud XR stack, encompassing both normal and impaired scenarios, to simulate real-world challenges.

In Chapter \ref{chap:ml_evaluation}, we leverage the collected cloud gaming datasets to explore the applicability of unsupervised ML models for detecting QoE degradations in cloud gaming applications. Unsupervised ML offers the advantage of not requiring labeled data, which is often challenging to obtain in practice. This chapter examines eight unsupervised ML models, focusing on three key criteria: detection accuracy, robustness to data contamination, and computational efficiency. By introducing the Window Anomaly Decision (WAD) method and evaluating model performance using metrics such as F1-score and Matthews Correlation Coefficient (MCC), the chapter provides insights and practical recommendations on selecting the most suitable unsupervised ML models based on the application and configuration requirements. The findings from this chapter are applied to a practical task within the MOSAICO project, not presented in this thesis: cloud gaming traffic detection. This task, essential for reducing latency and enhancing QoE in cloud gaming applications, is reformulated as an unsupervised learning problem. This approach addresses the generalization limitations of supervised methods by achieving greater accuracy to unseen or novel traffic patterns. This work has been published at an international conference \cite{ky2023hybrid}.

%Chapter \ref{chap:cg_detection} extends the findings from the previous chapter to address a practical task within the MOSAICO project: cloud gaming traffic detection. This task, essential for reducing latency and enhancing QoE in cloud gaming applications, is reformulated as an unsupervised learning problem. This approach addresses the limitations of supervised methods, which often fail to generalize to unseen or novel traffic patterns. The chapter also introduces a hybrid architecture that combines a P4-based hardware module for high-speed packet processing with a \ac{NFV} module, enabling the application of ML models for traffic classification. This architecture achieves high processing speed and detection accuracy, demonstrating its effectiveness in real-world scenarios.

To further improve anomaly detection (AD) in multivariate time series and address limitations such as limited robustness to data contamination, Chapter \ref{chap:cats} proposes a novel framework called CATS (Contrastive Anomaly detection on Time Series). CATS leverages contrastive learning to enhance time series representation learning. Its core innovation lies in a custom loss function that clusters temporally similar time series windows while pushing apart dissimilar points in the latent space. The framework also incorporates negative data augmentation techniques to generate synthetic anomalies. Evaluated on public benchmark datasets and the cloud gaming datasets collected in this thesis, CATS achieves superior performance compared to state-of-the-art AD models, while maintaining robustness to data contamination.

Finally, in Chapter \ref{chap:diagnosis} focuses on root cause diagnosis in Cloud VR applications operating on Wi-Fi networks. We propose a novel two-stage framework that integrates the CATS anomaly detection model from Chapter \ref{chap:cats} with a lightweight supervised classifier, such as a Support Vector Machine (SVM). The first stage detects anomalies in Cloud VR KPIs collected from Orange's Livebox, while the second stage classifies these anomalies into specific categories.  Evaluated on datasets collected under controlled Cloud VR experiments, this framework outperforms existing one-stage and two-stage approaches. Notably, it performs well even with limited labeled data, making it suitable for real-world deployment. Furthermore, its low inference and retraining times ensure practicality for dynamic network environments where new anomaly classes may emerge.

%In summary, this thesis provides a comprehensive suite of solutions to address the challenges of detecting anomalies and diagnosing root causes in low-latency applications. These contributions pave the way for more reliable network operations and enhanced user experiences in the era of high-demand, latency-sensitive services.




